\section{Сбор и подготовка данных}

Для оценки качества работы пайплайна и возможности
сравнения его с аналогами необходимо создать репрезентативный датасет,
который бы удовлетворял требованиям, описанным в главе 1.4.

\subsection{Поиск данных}

В качестве первичной гипотезы рассматривалось использование существующих открытых
пользовательских хранилищ в формате Obsidian с их последующей модификацией и фильтрацией.
В ходе поиска было отобрано несколько качественных репозиториев:
\begin{itemize}
      \item личные заметки общего назначения~\cite{noodleslove_notes};
      \item структурированная документация коллектива \textit{Hundred Rabbits}~\cite{100r_co};
      \item корпус кулинарных рецептов \textit{based.cooking}~\cite{based_cooking};
      \item заметки по программной инженерии~\cite{keyvanakbary_notes}.
\end{itemize}

После их детального анализа выявился один общий критический недостаток:
нарушалось правило ``заметка = сущность''.
Наиболее распространенной причиной его нарушения были
склейки текстов на разные тематики в одном файле.

Помимо структурного несоответствия был выявлен ряд дополнительных недостатков:
\begin{itemize}
      \item Семантическая разреженность данных.
            Некоторые хранилища состояли в основном из дневниковых записей, заметок или рецептов,
            между которыми практически не было связей, либо они были неоднозначными.
      \item Несоответствие объема.
            Почти все найденные хранилища были либо слишком малы, либо слишком велики.
            К репозиториям с недостающим объемом пришлось бы искусственно добавлять данные, а из
            слишком больших удалять,
            что потенциально может создать массу висячих ссылок и
            спровоцировать потерю ценных связей и контекста.
      \item В некоторых репозиториях отсутствует явная указанная лицензия
            либо присутствуют фрагменты текстов или кода,
            потенциально нарушающие авторские права.
            Это делает обработку этих данных рискованной.
\end{itemize}

\textbf{Выбор Википедии как основного источника данных.}
В качестве следующего источника данных было решено рассмотреть Википедию.
У нее есть ряд важных преимуществ:
\begin{itemize}
      \item Естественная атомарность статей. Зачастую статьи соответствуют одной сущности,
            что согласуется с правилом ``одна заметка = одна сущность''.
      \item Наличие ссылок. В статьях изначально есть гиперссылки на другие статьи.
            Эти ссылки образуют граф,
            к которому можно получить доступ и использовать его как для сбора данных,
            так и для первоначальной разметки.
      \item Качество содержимого. Подавляющая часть популярных статей
            написана в энциклопедическом стиле,
            обладает высокой плотностью информации, ее последовательным и связным изложением.
\end{itemize}

Все вышеперечисленное делает Википедию
подходящим источником данных для валидационного датасета.

\subsection{Создание валидационного датасета}

\textbf{Построение датасета и восстановление связей.}
Формирование базового датасета проводилось при помощи Wikipedia API,
обращение к которому производилось посредством библиотеки \texttt{wikipediaapi}.
Данный инструмент позволяет сохранять тексты статей, а также предоставляет доступ к графу из ссылок.
Важно отметить, что граф поддерживает типизацию ребер, что также поможет формированию датасета.
Однако лишь малая его часть размечена, и в основном это ссылки, соединяющие популярные статьи.

Важно отдельно отметить, что статьи имеют специальный раздел, содержащий краткую сводку (summary).
Он будет использован качестве контекста пайплайном GraphRAG
для фильтрации заметок на этапе сбора контекста.

\textbf{Формирование структуры элементов датасета.}
Поскольку ресурсы ограничены,
необходимо определить несколько наиболее показательных сценариев для оценки алгоритма.
На основе анализа особенностей работы GraphRAG, описанных в главе 1.2,
был составлен список наиболее показательных структур
элементов валидационного датасета
и их содержимого:

\begin{itemize}
      \item \textbf{Hub-n-spoke (Hub)}. Структура в форме граф в виде звезды,
            состоящая из статьи-хаба и ее соседей.
            Главная задача алгоритма в этом случае - идентификация хаба.
            Алгоритмы GraphRAG зачастую пытаются обнаружить именно такие узлы,
            так как они дают непосредственный доступ к большому количеству информации.
            Тематика датасета -- производство полупроводников и микроэлектроника.
            Итоговый набор статей включает как общие концепции (закон Мура, фотолитография, кремний),
            так и конкретные технологии или компании (ASML, TSMC, транзисторы, архитектура ARM, Apple M1).
      \item \textbf{Versus}. Структура, состоящая из двух разных тематик, которые местами могут показаться схожими.
            Например, Apple как компания и apple как яблоко.
            Такой подход заставит алгоритм работать с содержимым заметок, а не только с названиями сущностей.
            В собранном наборе данных эта структура реализована на примере явления омонимии:
            корпорация Apple (включая связанные узлы, например, Стив Джобс, iPhone, macOS)
            противопоставляется яблоку как фрукту (яблочный пирог, сорт Грэнни Смит)
            и связанным с ним историческим фактам (Исаак Ньютон).
      \item \textbf{Realistic}. Данный элемент будет содержать две схожие тематики.
            Его задача - проверка созданных алгоритмом связей на шум.
            Для датасета было решено связать
            концепции искусственного интеллекта (перцептроны, машинное обучение,
            обратное распространение ошибки и т.д.) с их историей.
      \item \textbf{Sequence}. Образец,
            состоящий из последовательных или вытекающих друг из друга событий или терминов.
            В качестве предметной области для этого набора была выбрана хронология развития представлений
            о строении атома (от открытия катодных лучей и электрона до пудинговой модели,
            экспериментов Резерфорда и модели Бора).
            Этот подкорпус будет нацелен на проверку способности алгоритма находить сложные, многошаговые связи.
      \item \textbf{Personal}. Элемент,
            созданный в ручную по аналогии с ежедневными заметками и содержащий различные события,
            планы, наброски идей.
            В нем смоделированы типичные личные записи: привязки к датам,
            повседневные задачи (починка велосипеда, покупки), встречи (с научным руководителем, знакомыми)
            и проекты (работа над дипломом, исследование малых LLM).
            Данный образец поможет оценить работу в условиях малого и специфичного контекста,
            отличного от текстов статей.
\end{itemize}

Визуализации ГЗ,
созданных на основе этих подкорпусов,
изображена на рисунке~\ref{fig:all_dataset_graphs}.

\subsubsection{Алгоритм генерации связей}

Для формирования датасета был написан скрипт для выгрузки и
обработки статей с использованием \texttt{wikipediaapi},
поддерживающий два режима сохранения:

\begin{itemize}
      \item выгрузка заданного набора статей
      \item обход соседей заданной статьи в ширину с контролем глубины обхода
\end{itemize}

Во второй стратегии сохранения
соседи статьи ранжируются с помощью Cross Encoder модели на основе их краткой сводки
и отбирается \textit{N} наиболее релевантных,
что предотвращает уход алгоритма от тематики корневой статьи.
Далее алгоритм увеличивает глубину рекурсии и аналогично обрабатывает каждого отобранного соседа
до тех пор пока либо не сохранит заданное число статей либо не обойдет все статьи до заданной глубины рекурсии.

Каждая статья, прошедшая проверку Cross Encoder подвергается обработке.
В нее входит:
\begin{itemize}
      \item Нормализация названий файлов, включающая замену пробелов на подчеркивания и
            удаление некоторых служебных символов для
            однозначного соответствия названий статей названиям файлов.
      \item Очистка текста от сносок и маркеров (например, \texttt{[1]}),
            а также удаление служебных секций статьи, которые не несут
            содержательной информации для GraphRAG
            (например, ``См. также'', ``Литература'', ``Ссылки'').
      \item Формирование специального блока ссылок, извлеченных из графа Википедии, на другие статьи.
            В данном блоке остаются ссылки только на сохраненные статьи,
            что решает проблему висячих ссылок.
\end{itemize}

После сохранения статей начинается этап их ручной доразметки.
Этот шаг включает:
\begin{itemize}
      \item Анализ и фильтрация ссылок, сохраненных из графа Википедии.
            Все ссылки, для идентификации которых нет контекста в содержимом статьи, удаляются.
            Например, в статье Билл Гейтс сказано, что он соосновал Microsoft, следовательно связь
            Билл Гейтс -> (соосновал) Microsoft должна быть оставлена.
      \item Анализ типов ссылок. Как было упомянуто ранее,
            в графе Википедии часть ссылок уже типизирована.
            Эти типы дополняются недостающими, а также абстрагируются для сокращения их общего числа.
            Критично обеспечить оптимальный уровень абстракции,
            чтобы не запутать языковая модель:
            большое количество узких по смыслу ссылок может показаться лучшим решением,
            но на практике модель начнет выбирать неверные типы, например из-за недостатка контекста.
\end{itemize}

Для упрощения разметки и алгоритмов
все типы отношений
фиксируются на английском языке
(например, \texttt{Uses}, \texttt{Part of}, \texttt{Mentions}),
несмотря на то, что содержимое
некоторых заметок на русском языке.

В каждом подкорпусе присутствует файл .meta,
в котором содержится список статей с Википедии,
на основе которых он был создан,
а также список допустимых типов отношений.

Результатом ручной доразметки является блок классифицированных ссылок
в каждой статье,
оформленный в синтаксисе плагина Dataview:

\begin{verbatim}
## Related Connections:
- Mentions:: [[ASML_Holding]]
- Manufactures:: [[Integrated_circuit]]
- Sustains:: [[Moores_law]]
- Is a:: [[Photolithography]]
\end{verbatim}

\begin{table}[H]
      \centering
      \begin{tabular}{lcccc}
            \hline
            Подкорпус              & \# заметок & \# ссылок & \# типов связей & Ср. исх. ссылок на заметку \\
            \hline
            \texttt{hub\_n\_spoke} & 15         & 77        & 12              & 5.13                       \\
            \texttt{realistic}     & 15         & 79        & 12              & 5.27                       \\
            \texttt{versus}        & 15         & 44        & 14              & 2.93                       \\
            \texttt{sequence}      & 7          & 23        & 7               & 3.29                       \\
            \texttt{personal}      & 11         & 23        & 7               & 2.09                       \\
            \hline
            \textbf{Итого}         & 63         & 246       & ---             & 3.90                       \\
            \hline
      \end{tabular}
      \caption{Сводная статистика по эталонных подкорпусах}
      \label{tab:dataset_samples}
\end{table}


В таблице~\ref{tab:dataset_samples} приведена информация о собранных данных.
\texttt{hub\_n\_spoke} и \texttt{realistic} обладают самым высоким средним числом связей на заметку
(в среднем 5.2 исходящих ссылки),
то есть в них много семантически близких и связных сущностей.
Они помогают оценить пайплайн построения графа на устойчивость к шуму.
\texttt{versus}, несмотря на идентичное число заметок, имеет вдвое меньше
связей, поскольку в нем часть сущностей не имеют никаких семантических отношений.
Аналогичная ситуация в \texttt{sequence}: в нем изложена последовательность,
в которой каждая сущность в основном связана только с предыдущей и следующей.

Датасеты хранятся в директории \texttt{data/test\_vaults/}.


\subsubsection{Алгоритм извлечения контекста}

Для оценки retrieval-этапа был создан отдельный датасет,
состоящий из фрагментов текстов, в которых явно или косвенно
упоминаются сущности и подкорпуса.

В качестве минимального фрагмента текста было решено взять предложение,
содержащие упоминание. Это позволит сохранить контекст,
а также проверять алгоритм на точность извлечения,
то есть долю других предложений, не содержащих упоминание.

Для каждого элемента датасета было проведено
ручное извлечение всех фрагментов с упоминаниями.
На их основе создан набор JSON-файлов
(по одному на элемент эталонного набора данных для оценки алгоритма генерации ссылок),
хранящийся в директории \texttt{data/test\_retrieval/gold/}.
Каждый файл содержит массив кортежей \((s, e, sent, is\_direct)\),
где \(s\) -- файл-источник,
\(e\) -- целевая сущность,
\(sent\) -- предложение, в котором упомянута сущность,
\(is\_direct\) -- содержит ли предложение прямое упоминание сущности.
Допускается, что для одной пары
\((s,e)\) может присутствовать несколько
записей с разными извлеченными предложениями.

Важно уточнить, что для соблюдения требований к данным,
описанным в главе 1.4,
в датасете используются только чанки,
содержащие упоминания сущностей, представленных в подкорпусе.

\subsubsection{Алгоритм LLM-классификации отношений}

Для валидации LLM-классификатора отношений используется дополненная
версия описанного выше датасета для оценки retrieval.
В каждую запись в JSON-файлах было добавлено поле
\texttt{link\_type}, соответствующее истинному типу отношения.
Таким образом финальная структура кортежей имеет следующий вид:
\((s, e, sent, is\_direct, link\_type)\).

Значения \(link\_type\) были взяты из блока ссылок
датасета для оценки алгоритма генерации связей.
В случае, если сущность упомянута в нескольких
предложениях, каждое из них получает одинаковое значение
\(link\_type\).

\textbf{Вывод.}

Для оценки качества и устойчивости алгоритмов к разным сценариям использования были сформированы
датасеты, различающиеся тематикой, плотностью информацией, сложностью и набором сущностей,
для оценки как всего пайплайна, так и для конкретных его этапов.
Данный подход задуман для упрощения тонкой настройки гиперпараметров и
выбора оптимальных стратегий извлечения контекста и классификации связей.

Выбор в пользу Википедии обеспечил выполнение правила ``одна заметка = одна сущность''
и позволил получить управляемый объем данных с качественным содержимым, достаточным для оценки GraphRAG.



% \textbf{Требования к данным.}
% С учетом устройства предлагаемого конвейера (см. раздел~2.1) к датасету были предъявлены следующие требования:
% \begin{enumerate}
%     \item \textit{Структурная согласованность.} Основная единица хранения и обработки в системе~--- Markdown-файл (``заметка'').
%           На уровне графа знаний вершина однозначно идентифицируется путем файла,
%           поэтому требуется соблюдение принципа ``одна заметка = одна сущность'' (атомарность представления).
%     \item \textit{Наличие смысловых связей.} Корпус должен содержать нетривиальные межтекстовые отношения
%           (иерархические, причинно-следственные, функциональные и др.), иначе оценка качества типизации
%           становится статистически неустойчивой и малоинформативной.
%     \item \textit{Контролируемый объем.} Для построения ``золотого стандарта''
%           необходима ручная или экспертно контролируемая разметка.
%           Следовательно, требуется ограниченный по размеру подграф,
%           в котором разметка выполнима в рамках дипломного исследования.
%     \item \textit{Качество и связность текста.} Языковые модели и модели ранжирования устойчивее работают на связных описательных фрагментах.
%           Преобладание обрывочных записей, фрагментов кода и цитат без контекста приводит к росту доли ложноположительных и ложноотрицательных связей.
%     \item \textit{Юридическая воспроизводимость.} Источник должен иметь ясный лицензионный статус,
%           допускающий обработку и публикацию производных данных (в частности, размеченных подмножеств и статистик).
% \end{enumerate}

% \textbf{Ключевое ограничение ``заметка = сущность''.}
% В предлагаемой архитектуре идентификатором сущности выступает файл (его относительный путь в хранилище).
% Это допущение используется системно: при индексировании и разбиении на чанки каждый фрагмент текста наследует идентификатор файла;
% при генерации кандидатных связей сопоставляются пары ``файл-источник~--- файл-цель'';
% при экспорте результаты записываются обратно в Markdown в виде типизированных ссылок вида \texttt{Relation:: [[TargetNote]]}.
% Если одна заметка содержит множество разнотипных сущностей (например, дневниковую запись с перечислением задач, людей и событий)
% или, наоборот, одна сущность ``распылена'' по нескольким файлам, то возникает неоднозначность:
% \begin{itemize}
%     \item невозможно однозначно интерпретировать ребро графа как отношение между сущностями (ребро связывает \textit{документы}, а не \textit{концепты});
%     \item при разметке и оценке качества предсказаний существенно возрастает субъективность: эксперт вынужден решать, к какой именно сущности внутри текста относится классифицируемая связь;
%     \item рефакторинг исходного корпуса (дробление лонгридов, переименование, перенос фрагментов, нормализация ссылок) фактически превращается в создание нового датасета и искажает ``естественную'' топологию связей, которую требуется моделировать.
% \end{itemize}
% Таким образом, нарушение принципа атомарности в исходном хранилище является не частным неудобством, а фундаментальным методологическим препятствием для построения валидируемого ``золотого стандарта''.

% \textbf{Предварительный анализ открытых Obsidian-хранилищ.}
% В качестве первичной гипотезы рассматривалось использование существующих пользовательских хранилищ (vaults) формата Obsidian, опубликованных в open-source. На этапе отбора наличие или отсутствие уже проставленных ссылок не считалось критичным: предполагалось, что связи можно извлечь и типизировать путем ручной или полуавтоматической разметки.

% В ходе поиска были отобраны и проанализированы несколько публичных репозиториев:
% \begin{itemize}
%     \item \url{https://github.com/noodleslove/notes/} (личные заметки общего назначения);
%     \item \url{https://github.com/hundredrabbits/100r.co} (структурированная документация коллектива);
%     \item \url{https://github.com/lukesmithxyz/based.cooking} (корпус кулинарных рецептов);
%     \item \url{https://github.com/keyvanakbary/learning-notes} (заметки по программной инженерии).
% \end{itemize}

% Детальный анализ показал, что такие хранилища систематически не удовлетворяют ключевому требованию атомарности ``заметка = сущность''. 
% Типичные паттерны нарушений включают: дневниковые записи с множеством сущностей в одном документе, ``оглавления'' и MOC-страницы (map-of-content), 
% лонгриды с несколькими связанными темами, а также файлы, преимущественно состоящие из фрагментов кода и выписок. 
% Адаптация подобных хранилищ потребовала бы радикальной переработки: разбиения документов на атомарные сущности, переименования, 
% восстановления и перенормализации ссылок, а иногда полного переписывания содержания, что противоречит цели исследования 
% (оценивать алгоритм на реалистичных данных, а не на искусственно подготовленном корпусе).

% Помимо структурного несоответствия были выявлены дополнительные ограничения:
% \begin{itemize}
%     \item \textit{Семантическая разреженность.} Во многих доменных корпусах (например, рецепты) набор возможных отношений крайне узок и не позволяет проверять устойчивость типизации для разнообразных классов связей.
%     \item \textit{Фрагментация при субсэмплинге.} Реальные пользовательские хранилища часто велики по объему; при выделении подмножества возникает высокий риск ``повисших'' ссылок (dangling links) и потери контекста, что смещает оценочные метрики.
%     \item \textit{Зашумленность текста.} Обрывочные формулировки, телеграфный стиль, вставки кода и цитаты без контекста усложняют извлечение сущностей и повышают неопределенность в решениях моделей.
%     \item \textit{Лицензирование.} Во многих персональных репозиториях отсутствует явная лицензия либо присутствуют фрагменты, потенциально нарушающие авторские права, что делает публикацию производных размеченных данных рискованной.
% \end{itemize}

% \textbf{Выбор Википедии как базового источника корпуса.}
% Для выполнения требований к атомарности и воспроизводимости 
% был выбран корпус статей Википедии, преобразованный в формат Obsidian-хранилища. 
% Такой выбор обоснован следующими свойствами:
% \begin{itemize}
%     \item \textit{Естественная атомарность.} Статья энциклопедии в большинстве случаев соответствует одной сущности/концепту, что согласуется с принципом ``одна заметка = одна сущность''.
%     \item \textit{Наличие базовой топологии.} Гиперссылки между статьями образуют связный граф, который удобен как источник кандидатных ребер и как слабая разметка для последующей валидации.
%     \item \textit{Качество текста.} Энциклопедический стиль обеспечивает высокую информационную плотность и связность изложения, что благоприятно для извлечения контекстов и принятия решений LLM.
%     \item \textit{Воспроизводимость сбора.} Сбор может быть формализован через публичный API, фиксируемые параметры обхода и детерминированные правила очистки/нормализации.
% \end{itemize}

% \textbf{Построение Markdown-корпуса и восстановление связей.}
% Подготовка корпуса выполнялась программно на основе Wikipedia API с использованием библиотеки \texttt{wikipediaapi} и включала следующие шаги.
% \begin{enumerate}
%     \item \textit{Выбор подграфа статей.} Реализованы два режима:
%           (i) целевая выгрузка заданного набора заголовков статей;
%           (ii) обход ``hub-and-spoke'' от одной стартовой темы (\texttt{seed}) с ограничением глубины и числа страниц.
%           Во втором режиме исходящие ссылки ранжируются кросс-энкодером (reranker), что позволяет сохранять более релевантные теме статьи при фиксированном бюджете на число заметок.
%     \item \textit{Нормализация идентификаторов.} Имена файлов приводятся к стабильному формату (замена пробелов на подчеркивания, удаление некоторых служебных символов), чтобы wikilink'и однозначно соответствовали путям файлов.
%     \item \textit{Очистка текста.} Из текста удаляются маркеры сносок (например, \texttt{[1]}), нормализуются подряд идущие переводы строк. Дополнительно из конца статей отсекаются секции, не несущие содержательной информации для извлечения отношений (например, ``См. также'', ``Литература'', ``Ссылки'' и их английские аналоги).
%           % \item \textit{Формирование блока связей.} Для каждой сохраненной статьи формируется секция \texttt{## Related Connections} со списком исходящих wikilink'ов. Для предотвращения ``повисших'' ребер ссылки на статьи, отсутствующие в выбранном подмножестве, удаляются.
%     \item \textit{Фильтрация ссылок по тексту (опционально).} Для снижения доли навигационных и слабосодержательных ссылок предусмотрен режим, в котором исходящая ссылка сохраняется только при наличии упоминания целевой сущности в основном тексте статьи (на основе лемматизации и проверки близости словоформ в предложениях).
% \end{enumerate}

% \textbf{Формирование ``золотого стандарта'' для оценки.}
% Для оценки качества системы требуются данные в двух представлениях.
% \begin{itemize}
%     \item \textit{Граф как Obsidian-хранилище.} Эталонные ребра фиксируются непосредственно в Markdown-файлах в виде типизированных строк внутри секции связей: \texttt{- relation:: [[Target]]}. Такая форма разметки позволяет (i) сохранять данные в естественном для Obsidian виде и (ii) напрямую сравнивать предсказания системы с эталоном путем парсинга Markdown.
%     \item \textit{Датасет контекстов упоминаний для retrieval-оценки.} Для каждой связи ``источник~--- цель'' извлекаются предложения контекста, в которых в тексте источника упоминается целевая сущность. Данные сериализуются в JSON и включают поля \texttt{source\_file}, \texttt{entity}, \texttt{sentence}, \texttt{is\_direct\_mention} и (после добавления типизации) \texttt{link\_type}.
% \end{itemize}
% Такое раздельное представление позволяет оценивать как задачу \textit{link classification} (типизация ребер), так и качество выделения релевантных контекстов, необходимых для LLM-классификатора и для GraphRAG-поиска.

% \textbf{Датасет и протокол оценки retrieval.}
% Оценка retrieval-компонента формализуется как задача воспроизведения контекстов упоминаний: для каждой пары \((s, e)\), где \(s\)~--- файл-источник (\texttt{source\_file}), а \(e\)~--- целевая сущность (\texttt{entity}), система должна вернуть хотя бы одно предложение/фрагмент текста (\texttt{sentence}), семантически совпадающий с эталонным контекстом из ``золотого стандарта''.
% Эталонные данные хранятся в виде набора JSON-файлов (по одному на подкорпус) в директории \texttt{data/test\_retrieval/gold/}. Каждый элемент массива соответствует одной ``истинной'' паре \((s, e)\) и содержит конкретное предложение из текста источника, в котором упоминается сущность \(e\).

% Предсказания retrieval-стратегий вычисляются непосредственно на Markdown-хранилищах подкорпусов (директория \texttt{data/test\_vaults/gold/}) и сохраняются в том же JSON-формате. В экспериментальном протоколе рассматриваются три стратегии:
% \begin{itemize}
%     \item \textit{Strict}~--- извлечение кандидатных упоминаний на основе лемматизированного совпадения сущностей в тексте;
%     \item \textit{Broad}~--- извлечение кандидатных упоминаний через векторный поиск по чанкам с последующим переранжированием кросс-энкодером;
%     \item \textit{Combined}~--- объединение \textit{Strict} и \textit{Broad} с устранением дубликатов.
% \end{itemize}

% Сопоставление эталона и предсказаний производится не по точному совпадению строк, а по ``нечеткому'' критерию на уровне лемм. Для каждого эталонного предложения вычисляется множество лемм (без стоп-слов и пунктуации) и проверяется, что доля совпавших лемм между эталоном и предсказанным фрагментом превышает фиксированный порог. Такая процедура делает метрику устойчивой к незначительным различиям в пунктуации и форматировании и к тому, что retrieval может возвращать фрагмент, покрывающий эталонное предложение не дословно.
% Итоговая оценка агрегируется через микро-метрики \(Precision\), \(Recall\) и \(F1\) по счетчикам \(TP/FP/FN\), где:
% \begin{itemize}
%     \item \(TP\): эталонный контекст для пары \((s,e)\) найден (существует предсказанный фрагмент, прошедший ``нечеткое'' сопоставление);
%     \item \(FN\): эталонный контекст отсутствует среди предсказаний;
%     \item \(FP\): предсказана пара \((s,e)\), для которой не удалось сопоставить ни один предсказанный фрагмент с эталонными предложениями.
% \end{itemize}

% \textbf{Датасет и протокол оценки LLM-классификации отношений.}
% Оценка LLM-компонента трактуется как многоклассовая классификация типа отношения для пар \((s,e)\) при условии, что retrieval уже предоставил набор релевантных контекстов. В качестве ``золотого стандарта'' используется тот же набор JSON-файлов \texttt{data/test\_retrieval/gold/}, но дополненный меткой \texttt{link\_type}, которая соответствует типизированной ссылке в Markdown-хранилище и интерпретируется как истинный класс (\texttt{true\_type}).

% Для каждой пары \((s,e)\) все доступные предложения-контексты агрегируются и далее подаются в модуль ранжирования контекстов. Ранжирование выполняется кросс-энкодером (модель переранжирования из конфигурации проекта), который получает на вход запрос вида ``What is the relationship between \dots?'' и набор кандидатных предложений; после чего выбираются \textit{top-n} наиболее релевантных предложений. При ранжировании дополнительно приоритизируются предложения, отмеченные признаком прямого упоминания (\texttt{is\_direct\_mention}).

% Существенной особенностью протокола является то, что пространство классов задается не глобально, а на уровне каждого подкорпуса. Для каждого sample список допустимых типов связей \texttt{llm\_link\_types} извлекается из файла конфигурации \texttt{.kg\_builder/config.json}, что позволяет:
% \begin{itemize}
%     \item ограничивать выбор LLM только релевантными типами для конкретной предметной области;
%     \item исключать классы, отсутствующие в эталонной разметке данного поднабора;
%     \item корректно сравнивать результаты между подкорпусами, не смешивая различные онтологии связей.
% \end{itemize}

% Запросы к LLM выполняются батчами; после каждого батча результаты сохраняются в JSONL-файл-чекпоинт, что обеспечивает воспроизводимость и возможность продолжить оценку после прерывания. Для качества классификации вычисляются precision/recall/F1 по каждому классу, а также macro- и weighted-усреднения (веса пропорциональны поддержке класса в эталоне). Метрики рассчитываются отдельно по каждому подкорпусу и в совокупности по всем подкорпусам.

% \textbf{Контрольные подкорпуса.}
% Для проверки устойчивости алгоритма к различным видам структур и сложностей были сформированы несколько небольших ``контрольных'' хранилищ, каждое из которых соответствует отдельному поднабору (sample). Подкорпуса различаются тематикой, плотностью ссылок и специфическими ``сложными случаями'' (например, омонимия и конкурирующие смыслы сущностей).
% Сводная статистика по использованным подкорпусам приведена в таблице~\ref{tab:dataset_samples}.

% \begin{table}[H]
%     \centering
%     \caption{Сводная статистика контрольных хранилищ (``золотой стандарт'')}
%     \label{tab:dataset_samples}
%     \begin{tabular}{lccc}
%         \hline
%         Подкорпус              & Число заметок & Число типизированных ссылок & Число контекстных пар \\
%         \hline
%         \texttt{hub\_n\_spoke} & 15            & 77                          & 96                    \\
%         \texttt{realistic}     & 15            & 79                          & 103                   \\
%         \texttt{versus}        & 15            & 44                          & 62                    \\
%         \texttt{sequence}      & 7             & 23                          & 44                    \\
%         \texttt{personal}      & 11            & 23                          & 22                    \\
%         \hline
%     \end{tabular}
% \end{table}

% Итоговый выбор в пользу Википедии и построенных на ее основе контрольных хранилищ обеспечил выполнение ключевого методологического требования атомарности сущностей и позволил получить воспроизводимый набор данных с управляемым объемом разметки, достаточный для количественной оценки качества GraphRAG-конвейера.
